%% kapa: exam=1 round=23 year=2012 period=2 form=A subject=accounting subject_ko=회계학 seq=21 number=21
%% source: Q-Net 감정평가사 기출문제 (cst00309) artlSeq=5052580 file=raw/1008/5052580/2012년 제23회 감정평가사 제1차시험 기출문제(2교시 A형).pdf
%% conversion: kordoc 4.24.51 -> tools/gen1.py (auto); review=proofread-sonnet; flags=none
\begin{question}{21}
(주)한국은 20×6년 초 (주)서울을 매수하기로 하였다. 다음 자료를 이용하여 초과이익을 고려한 (주)한국의 (주)서울에 대한 매수금액은?
\begin{framedtext}
ㆍ(주)서울의 20×5년 말 자산의 공정가치 합계는 ￦1,500,000이고, 부채의 공정가치 합계는 ￦200,000이다. \par
ㆍ과거 5년간 (주)서울의 당기순이익은 아래와 같으며, 20×3년 당기순이익에는 중단사업이익 ￦200,000이 포함되어 있다. \par
\begin{center}\small
\begin{tabular}{@{}lrlr@{}}
\hline
ㆍ20×1년 & ￦150,000 & ㆍ20×4년 & ￦240,000 \\
ㆍ20×2년 & ￦100,000 & ㆍ20×5년 & ￦120,000 \\
ㆍ20×3년 & ￦300,000 & & \\
\hline
\end{tabular}
\end{center}
ㆍ정상이익률은 동일업종의 평균이익률인 10\%, (주)한국이 기대하는 투자수익률은 20\%이다. \par
ㆍ초과이익은 영구적이라고 가정한다. \par
\end{framedtext}
\begin{choices}
  \item ￦710,000
  \item ￦1,100,000
  \item ￦1,148,000
  \item ￦1,360,000
  \item ￦1,560,000
\end{choices}
\end{question}
\answer{4}
